\runninghead{Each Text in Its Setting}
\section{Each Text in Its Setting}
\label{sec:each-element}

Each appointed text below is set in the chapter or psalm it comes from, with
what the Missal keeps and leaves out and what the Fathers and later
commentators say of its words.

\subsection{Introit: Jeremias writes to the exiles}
\label{sec:element-introit}

The antiphon comes from a letter. Jeremias sent it from Jerusalem ``to the
residue of the ancients that were carried into captivity, and to the priests,
and to the prophets, and to all the people, whom Nabuchodonosor had carried
away from Jerusalem to Babylon,'' after King Jechonias and the princes and
craftsmen of Juda had been taken there (Jer 29:1--2). The letter tells the
exiles to settle: ``Build ye houses, and dwell in them: and plant orchards,
and eat the fruit of them'' (v.~5). It tells them to marry and to have
children there, and to ``seek the peace of the city, to which I have caused
you to be carried away captives; and pray to the Lord for it: for in the peace
thereof shall be your peace'' (v.~7). It warns them against the prophets and
diviners among them, who prophesy falsely in God's name (vv.~8--9). Only then
come the promise of return and the words the Introit sings (vv.~10--14).

St Jerome reads verse~7 tropologically, of the Church in this world:
\latin{Quaeramus quoque pacem Ecclesiae \dots\ orare praecepit pro inimicis
nostris \dots\ ut non simus nostra tantum salute contenti, sed inimicorum
quoque quaeramus salutem}, ``let us seek the peace of the Church too \dots\ he
has commanded us to pray for our enemies \dots\ so that we may not be content
with our own salvation alone, but seek our enemies' salvation as well''
(\work{In Hieremiam} V, on 29:5--9, PL~24 cols.~891--892).

The antiphon keeps the three things God says he will do: think thoughts of
peace, hear those who call, and bring back the captivity. It begins
\latin{Dicit Dóminus}, ``The Lord saith,'' and so the whole antiphon is God
speaking in the first person, as it is in the letter. It omits the end and
patience, the going and praying, the seeking with the whole heart, the words
``I will be found by you,'' and the gathering from all nations.

The verse, Ps 84:2, belongs to a psalm ``for the sons of Core'' (84:1) that
begins with what God has done, ``Lord, thou hast blessed thy land: thou hast
turned away the captivity of Jacob. Thou hast forgiven the iniquity of thy
people: thou hast covered all their sins'' (vv.~2--3), and then prays for
it: ``Convert us, O God our saviour'' (v.~5). St Augustine, Cassiodorus and
the exposition printed under St Athanasius's name read the captivity as the
captivity of sin.

Blessed Ildefonso Schuster hears in the antiphon ``the divine answer to our
humble supplications'' and says that the Lord, ``even when he punishes, does
not forget his mercy, indeed he chastises because he loves'' (\work{The
Sacramentary} III, p.~186).

\subsection{Collect: a family that leans on grace alone}
\label{sec:element-collect}

The Collect has one petition, \latin{Famíliam tuam \dots\ contínua pietáte
custódi}, ``keep thy family with continual mercy,'' and one purpose, that the
family \latin{tua semper protectióne muniátur}, ``may always be defended by
thy protection.'' Between them stands a relative clause that gives the
family's condition: \latin{quæ in sola spe grátiæ cæléstis innítitur}, it
leans on the hope of heavenly grace, and on that alone. \latin{Innitor} is to
lean one's weight on something. \latin{Pietas} said of God is his
fatherly kindness; the Lasance renders it ``mercy.'' \latin{Contínua} and
\latin{semper} ask for protection without interruption.

\subsection{Epistle: what the elect of God put on}
\label{sec:element-epistle}

The third chapter of Colossians begins, ``Therefore if you be risen with
Christ, seek the things that are above'' (3:1). St Paul then tells the
Colossians to put to death what is earthly in them and to put away ``anger,
indignation, malice, blasphemy, filthy speech,'' stripping off ``the old man
with his deeds, and putting on the new'' (3:5--10). The Epistle begins where
he names what the new man puts on. The Missal opens it \latin{Fratres} and
drops the \latin{ergo} that ties verse~12 to what precedes.

St John Chrysostom comments on these verses in his eighth and ninth homilies
on Colossians. Paul calls the offence a \latin{querela}, a ``complaint,''
and Chrysostom says that the word makes it ``a petty matter.'' But by adding
``even as Christ forgave you'' Paul ``has persuaded us that even if we had
serious charges to bring, we ought to forgive,'' and not only to forgive but
``to love'' (\work{Hom.\ in Col.}\ 8). Charity is ``the bond of perfection''
because without it the other virtues ``fall asunder'': ``it is as in a ship,
even though her rigging be large, yet if there be no girding ropes, it is of
no service.''

At verse~15 the Missal and the Vulgate read \latin{pax Christi exsúltet in
córdibus vestris}, ``let the peace of Christ rejoice in your hearts.'' The
Greek text that Chrysostom expounds has a verb that means ``let it act as
umpire,'' and his commentary, with those of Theodoret and Theophylact who
follow him, is built on that word. St Thomas Aquinas expounds the Latin
\latin{exsúltet}, of the joy that follows from peace.

Verse~16 asks that the word of Christ dwell in the Colossians
\latin{abundánter}, and that they teach and admonish one another in psalms,
hymns and spiritual canticles, singing \latin{in córdibus}, in their hearts.
Verse~17 puts everything said or done under the name of the Lord Jesus. The
Missal ends the lesson \latin{per Iesum Christum Dóminum nostrum}, where the
Vulgate has \latin{per ipsum}, ``by him.''

\subsection{Gradual: thanks taken from a lament}
\label{sec:element-gradual}

Psalm~43 is a psalm ``for the sons of Core, to give understanding'' (43:1).
It begins with what the fathers have told: God planted his people in the land,
and ``they got not the possession of the land by their own sword,'' but by
his right hand (vv.~2--4). The speakers say they will not trust in their own
bow or sword (v.~7), and then come the two verses the Gradual sings: ``But
thou hast saved us from them that afflict us: and hast put them to shame that
hate us. In God shall we glory all the day long: and in thy name we will give
praise for ever'' (vv.~8--9). The next verse turns: ``But now thou hast cast
us off, and put us to shame: and thou, O God, wilt not go out with our
armies'' (v.~10). The lament runs to the end of the psalm and includes ``for
thy sake we are killed all the day long: we are counted as sheep for the
slaughter'' (v.~22), which St Paul quotes of the Christians in Rom 8:36.

The Gradual keeps the two verses of thanks and stops before the lament. It
reads \latin{Liberásti nos \dots\ ex affligéntibus nos} where the Vulgate has
\latin{salvasti enim nos de affligentibus nos}, the same reading that St
Ambrose expounds (\work{Enarr.\ in Ps.}\ 43.21, PL~14 col.~1154) and that
the psalm text printed with Cassiodorus has. Blessed Ildefonso Schuster calls the Gradual
``well adapted to the celebration of the victories of the martyrs''
(\work{The Sacramentary} III, p.~187).

\subsection{Alleluia: the cry from the depths}
\label{sec:element-alleluia}

Psalm~129 is one of the gradual psalms, \latin{Canticum graduum}, and the
sixth of the seven penitential psalms. The verse the Alleluia sings is its
opening: ``Out of the depths I have cried to thee, O Lord: Lord, hear my
voice.'' The Missal sings \latin{oratiónem meam}, ``my prayer,'' where the
Vulgate has \latin{vocem meam}. The psalm goes on, ``If thou, O Lord, wilt
mark iniquities: Lord, who shall stand it. For with thee there is merciful
forgiveness'' (vv.~3--4), and the Latin of that forgiveness is
\latin{propitiatio}.

St Augustine begins with Jonas, who ``cried from the deep; from the whale's
belly,'' and whose prayer the waves and the beast could not stop, ``since the
ears of God were in the heart of him who prayed.'' Then he names the deep:
``this mortal life is our deep.'' Some are so far down that they ``do not even
cry from the deep'' (\work{Enarr.\ in Ps.}\ 129.1). St Robert Bellarmine
compares the psalmist to a man calling from a deep well to someone on a high
mountain, and says that he cries from ``two depths, the depth of misery and
the depth of his heart'' (\work{Explanatio in Ps.}\ 129, vv.~1--2).

Blessed Ildefonso Schuster addresses God as the one ``who are drawn to us the
more powerfully, according as our misery is greater'' (\work{The
Sacramentary} III, p.~187).

\subsection{Gospel: the cockle among the wheat}
\label{sec:element-gospel}

The parable belongs to the day on which Jesus, sitting in a boat, taught the
crowds on the shore in parables (Mt 13:1--3). It comes after the sower and
before the mustard seed and the leaven (vv.~31--33). Jesus then sends the
crowds away and goes into the house, where the disciples ask him to explain
``the parable of the cockle of the field'' (v.~36). He answers: ``He that
soweth the good seed is the Son of man. And the field is the world. And the
good seed are the children of the kingdom. And the cockle are the children of
the wicked one. And the enemy that sowed them, is the devil. But the harvest
is the end of the world. And the reapers are the angels'' (vv.~37--39). The
angels ``shall gather out of his kingdom all scandals, and them that work
iniquity,'' and then ``shall the just shine as the sun, in the kingdom of
their Father'' (vv.~41, 43). None of this explanation is read in the Mass,
but it is the Lord's own, and the Fathers read the parable with it.

The Latin \latin{zizania} is the weed the Douay calls ``cockle'' and the
English translators of the Fathers call ``tares.'' St Jerome names it
\latin{lolium}, darnel, and says that while wheat and darnel are still in the
blade they are almost impossible to tell apart (\work{In Matthaeum} II, on
13:36--43, PL~26 col.~94). St Augustine says the same of every weed that
``spring[s] up in the manner of wheat, for instance darnel, for instance wild
oats, and all such as have the first leaf the same'' (\work{Enarr.\ in Ps.}\
64.15). The parable's own detail agrees: the cockle ``appeared'' only when the
blade had sprung up and borne fruit.

The Fathers name the cockle with different breadth. St John Chrysostom calls
it ``the societies of the heretics,'' sown by the devil beside the truth and
``in appearance somewhat like wheat'' (\work{Hom.\ in Mt.}\ 46.1), and St
Jerome reads the enemy's seed as \latin{haereticorum dogmata}, the teachings
of heretics. St Augustine, in his sermon on this Gospel, calls the tares
``evil Christians among the good'' and says that they are the same people as
the wayside, the stony ground and the thorns of the sower (\work{Sermo}~73.1--2);
in his questions on the Gospels he says that \latin{omnis immunditia in
segete, zizania dicuntur}, all the uncleanness in the crop is called cockle
(\work{Quaest.\ Ev.}\ I.10). St Thomas Aquinas calls them \latin{filii
nequam}, the wicked children, \latin{specialiter haeretici} (\work{Super Mt.}\
c.~13). All agree that the cockle looks like wheat and grows in the same field.

At 13:24--30 St Jerome declines to explain the parable before the Lord does:
\latin{Non ergo debemus praepropero intelligendi desiderio ante ejus notitiam
quaerere, quam a Domino disseratur}, ``we ought not, from a hasty desire to
understand, to seek its meaning before the Lord expounds it'' (PL~26
col.~89). St Hilary of Poitiers passes over the parables that the Lord himself
explained: \latin{De parabolis autem jam a Domino absolutis loqui otiosum
est} (\work{In Matthaeum} 13.1, PL~9 col.~993).

\subsection{Offertory: the same cry, sung again}
\label{sec:element-offertory}

The Offertory takes the same two half-verses of Psalm~129 as the Alleluia and
repeats the first: \latin{De profúndis clamávi ad te, Dómine: Dómine, exáudi
oratiónem meam: de profúndis clamávi ad te, Dómine}. With the Alleluia before
the Gospel, the words \latin{de profúndis clamávi} are heard three times
before the Secret. Blessed Ildefonso Schuster says that the gifts ``must be
accompanied by the incense of humility'' (\work{The Sacramentary} III,
p.~188). Rupert of Deutz and William Durandus hear in the Offertory the plea
for pardon that Joseph's brothers made to Joseph.

\subsection{Secret: offerings of propitiation}
\label{sec:element-secret}

The Secret names what is offered, \latin{hóstias \dots\ placatiónis},
``offerings of propitiation,'' and asks two things: that God, having mercy,
\latin{delícta nostra \dots\ absólvas}, ``absolve our sins,'' and
\latin{nutántia corda tu dírigas}, ``that thou direct our wavering hearts.''
\latin{Nutare} is to nod or sway, as a thing does that is not firmly set.
The first petition concerns what has been done; the second concerns the
heart that will act next. \latin{Placatio} is not the word of Ps 129:4,
\latin{propitiatio}, though both name what turns away God's anger.

\subsection{Communion: prayer that receives}
\label{sec:element-communion}

The Communion comes from the morning after Jesus cursed the barren fig tree.
Passing it, the disciples see it ``dried up from the roots,'' and Peter
remarks on it (Mk 11:20--21). Jesus answers: ``Have the faith of God,'' and
says that whoever tells a mountain to be cast into the sea ``and shall not
stagger in his heart, but believe,'' will have it done; ``therefore'' whatever
they ask in prayer, believing, they will receive (vv.~22--24). The next verse
adds: ``And when you shall stand to pray, forgive, if you have aught against
any man: that your Father also, who is in heaven, may forgive you your sins''
(v.~25). The antiphon sings verse~24 only, with \latin{quidquid} for the
Vulgate's \latin{omnia quaecumque}, and opens with the \latin{Amen dico
vobis} of verse~23.

St Bede reads verses 22--25 together and distinguishes those who pray:
\latin{Notanda distinctio deprecantium}. The perfect in faith can move
mountains. Those who are not yet perfect should ask that their sins be
forgiven them, \latin{et absque dubio impetrabunt quod petunt, si tamen in se
peccantibus aliis primo ipsi dimittunt}, ``and without doubt they will obtain
what they ask, provided that they first forgive others who sin against them''
(\work{In Marci evangelium expositio} III, on 11:22--25). St John
Chrysostom, on the same scene in Matthew, says that the fig tree was withered
``that He might train them to feel confidence,'' and that they would do far
more, ``being confident in faith and prayer'' (\work{Hom.\ in Mt.}\ 67.1--2).

\subsection{Postcommunion: the pledge and the effect}
\label{sec:element-postcommunion}

The Postcommunion distinguishes two things. What has been received
\latin{per hæc mystéria} is a \latin{pignus}, a pledge; what is asked for is
the \latin{efféctus}, the effect, of the salvation it pledges. A pledge is
given in advance of what it guarantees, and the prayer asks God for the
effect.
